\section{Metodologia}
\label{sc:metodologia}

O estudo é quantitativo e experimental. O desenho segue o de Şakar e Emekci~\cite{sakar2025rag}, com os mesmos seis métodos de RAG, o mesmo tamanho de fragmento e dois dos seus três conjuntos de dados, e acrescenta um controle sem recuperação, sete LLMs, um conjunto de dados em português e métricas comparadas diretamente com o gabarito. O modelo de embedding, o tamanho dos fragmentos e o número de fragmentos recuperados foram mantidos fixos em todos os experimentos, para isolar o efeito do método e do LLM.

\subsection{Conjuntos de dados}

A Tabela~\ref{tab:datasets} resume os três conjuntos. O critério comum foi que as perguntas e as respostas de referência viessem de um conjunto público, construído e revisado por terceiros, e que o corpus de recuperação fosse o próprio material de onde as respostas saem.

\begin{table}[!t]
\centering
\footnotesize
\setlength{\tabcolsep}{2.5pt}
\caption{Conjuntos de dados. Frag.: fragmentos de 1.000 caracteres no índice. Perguntas: total no conjunto e, entre parênteses, a amostra da camada principal. Gabarito: resposta escrita com os valores pedidos (Docugami), alternativa correta oficial (MedQA) e trecho do parágrafo (FaQuAD).}
\label{tab:datasets}
\begin{tabular}{lllrr}
\toprule
Conjunto & Domínio & Corpus & Frag. & Perguntas \\
\midrule
Docugami & financeiro (EN) & 20 relatórios 10-Q & 4.463 & 195 (100) \\
MedQA & médico (EN) & 18 livros-texto & 121.553 & 1.273 (100) \\
FaQuAD & educacional (PT) & 249 parágrafos & 318 & 900 (100) \\
\bottomrule
\end{tabular}
\end{table}

\textbf{Docugami SEC 10-Q.} O conjunto, distribuído pelo LlamaIndex \cite{docugami2024kgrag} e usado pelo estudo de referência, reúne 195 perguntas sobre os relatórios trimestrais (formulário 10-Q da SEC) de Apple, Amazon, Intel, Microsoft e NVIDIA, do terceiro trimestre de 2022 ao terceiro de 2023. As perguntas foram escritas por pessoas, e as respostas de referência foram geradas pelo GPT-4 e revisadas por pessoas. A maior parte pede valores das demonstrações financeiras, muitas comparando empresas ou trimestres, e cada resposta de referência cita o relatório de onde sai, o que permite medir a recuperação. Os PDFs foram convertidos em texto com o \texttt{pdftotext} preservando o leiaute, para que cada linha das tabelas financeiras mantivesse o rótulo junto dos valores.

\textbf{MedQA.} São questões reais do USMLE, o exame de licenciamento médico dos Estados Unidos, da divisão de teste do MedQA \cite{jin2021medqa}, na versão de quatro alternativas, com a alternativa correta oficial. O corpus de recuperação são os 18 livros-texto de medicina que acompanham o conjunto, na versão distribuída pelo MedRAG \cite{xiong2024mirage}, remontados livro a livro e refragmentados com o mesmo tamanho dos demais conjuntos. A pergunta enviada ao sistema é o caso clínico seguido das alternativas. O gabarito não indica de onde, nos livros, sai a resposta.

\textbf{FaQuAD.} O FaQuAD \cite{sayama2019faquad} é um conjunto de compreensão de leitura no formato do SQuAD \cite{rajpurkar2016squad}, com 900 perguntas escritas por pessoas sobre 249 parágrafos originais em português, extraídos de 18 documentos oficiais de uma universidade federal e de 21 artigos da Wikipédia sobre o ensino superior brasileiro. Cada resposta de referência é um trecho do parágrafo, e algumas perguntas têm mais de uma resposta aceita. Os parágrafos formam a base de recuperação, e o parágrafo de cada pergunta é o trecho certo para medir a recuperação.

O terceiro conjunto do estudo de referência, sobre o artigo do Llama~2, foi descartado porque tanto as perguntas quanto as respostas de referência foram geradas pelo GPT-4. O FaQuAD foi preferido a outros conjuntos em português, como o Pirá, por ter documentos originalmente escritos em português, e não traduzidos.

\textbf{Amostragem.} O desenho tem duas camadas. Na \emph{camada principal}, cinco LLMs de faixa de preço econômica respondem a 100 perguntas de cada conjunto, sorteadas com semente fixa. Na \emph{camada de generalização}, os dois LLMs intermediários respondem às 25 primeiras perguntas da mesma permutação, um subconjunto das 100, de modo que os sete modelos compartilham essas 75 perguntas e a comparação entre eles é pareada. A camada principal sustenta as conclusões sobre os métodos; a de generalização verifica se elas se mantêm em modelos mais capazes. No Docugami, 66 das 100 perguntas (e 14 das 25) têm algum valor numérico na resposta de referência; as demais são qualitativas e ficam fora da métrica principal desse conjunto.

\subsection{Pipeline de RAG}

Os documentos foram divididos com o \texttt{Recursive\-Character\-Text\-Splitter} em fragmentos de 1.000 caracteres com sobreposição de 100, os mesmos valores do estudo de referência, e codificados com o BGE-M3 \cite{chen2024bgem3}, modelo de embedding multilíngue que atende aos dois idiomas, usando apenas os vetores densos. Os vetores foram indexados no ChromaDB com distância euclidiana e o índice padrão do banco, o HNSW \cite{malkov2020hnsw}, de busca aproximada. Em todos os métodos, a recuperação devolve os $k = 4$ fragmentos mais próximos da pergunta.

Os seis métodos foram implementados diretamente sobre as APIs, sem \emph{frameworks} de orquestração, para que cada chamada ao LLM fosse contabilizada:

\begin{itemize}
  \item \textbf{Stuff} (1 chamada): os $k$ fragmentos são concatenados em um único prompt com a pergunta.
  \item \textbf{Refine} ($k$ chamadas): gera uma resposta com o primeiro fragmento e a refina com cada fragmento seguinte, alterando-a apenas se o novo fragmento trouxer informação relevante.
  \item \textbf{Map-Reduce} ($k+1$ chamadas): gera uma resposta parcial por fragmento e combina as parciais em uma resposta final.
  \item \textbf{Map-Rerank} ($k$ chamadas): gera, para cada fragmento isoladamente, uma resposta e uma nota de confiança de 0 a 100, e mantém a resposta de maior confiança.
  \item \textbf{Query Step-Down} (6 chamadas): a partir da pergunta e de um contexto inicial, gera quatro perguntas alternativas mais específicas, responde cada uma com o Stuff, com recuperação própria, e combina as quatro respostas em uma chamada final, que recebe a pergunta original e as respostas.
  \item \textbf{Reciprocal RAG} (5 chamadas): gera as mesmas quatro perguntas alternativas, responde cada uma com nota de confiança e mantém a resposta mais confiável, sem etapa de combinação.
\end{itemize}

O \textbf{controle sem RAG} (1 chamada) envia ao LLM apenas a pergunta, sem nenhum fragmento, com a mesma instrução do conjunto.

\textbf{Prompts.} Os prompts foram escritos em português e traduzidos literalmente para o inglês, idioma usado no Docugami e no MedQA: em um piloto com os prompts em português, parte dos modelos respondia em português às perguntas em inglês, o que penalizaria as métricas sem relação com o método. No Docugami e no FaQuAD, o modelo é instruído a responder usando apenas o contexto e a dizer que não sabe quando a resposta não estiver nele; no controle sem RAG, a dizer que não sabe quando não souber. No MedQA, como os livros-texto raramente trazem a resposta exata de um caso clínico, a instrução de recusa faria a acurácia medir recusas, e não conhecimento. Por isso, como no MIRAGE \cite{xiong2024mirage}, o modelo usa os trechos como apoio, recorre ao próprio conhecimento quando eles não bastam e sempre escolhe uma alternativa, começando a resposta pela letra, seguida de uma justificativa curta.

\subsection{Modelos de linguagem}

A Tabela~\ref{tab:llms} lista os sete LLMs. As faixas são de preço, e não de tamanho: o DeepSeek V4 Pro, um dos maiores modelos de pesos abertos, custa como um modelo econômico. A escolha buscou a diversidade necessária para generalizar os resultados, com uma justificativa para cada modelo:

\tabela{tab_llms}

\begin{itemize}
  \item \textbf{gpt-4o-mini}: um dos três LLMs do estudo de referência, o que liga os resultados aos dele.
  \item \textbf{GPT-6 Luna} e \textbf{Claude Haiku 5.5}: os modelos econômicos atuais das duas principais empresas de modelos fechados.
  \item \textbf{Qwen3.5-9B}: modelo pequeno, de pesos abertos (Apache-2.0), que pode ser executado localmente; garante que parte dos resultados seja reprodutível independentemente de uma API comercial \cite{palmer2024proprietary, chen2024chatgptbehavior}.
  \item \textbf{DeepSeek V4 Pro}: modelo de pesos abertos (MIT) com 1,6 trilhão de parâmetros em arquitetura de mistura de especialistas, a preço de modelo econômico.
  \item \textbf{Claude Sonnet 5} e \textbf{GPT-6 Sol}: os intermediários das mesmas duas famílias fechadas, o que permite medir o efeito do porte dentro de cada família, e testar se um modelo econômico com o método adequado alcança o intermediário \cite{chen2023frugalgpt}.
\end{itemize}

Para a camada de generalização, os intermediários foram preferidos aos modelos de topo das mesmas famílias (Claude Opus 4.8 e GPT-6 Astra), que custam de 2,5 a 5 vezes mais por token e, no caso do Astra, não permitem desligar o raciocínio. Pelo mesmo motivo, foram usados o Sonnet 5 e o GPT-6 Sol, e não as versões mais novas (Sonnet 5.5 e GPT-6.1 Sol), que também não permitem desligá-lo.\todo{Confirmar com o orientador a troca do Opus 4.8 e do GPT-6 Astra pelos intermediários.}

Em todos os modelos, o raciocínio estendido foi desligado, para comparar os métodos de RAG, e não o quanto cada modelo ``pensa'' antes de responder. A temperatura foi fixada em 0 nos modelos que aceitam o parâmetro; os modelos Claude e GPT-6 rejeitam a temperatura 0 e rodaram com a temperatura padrão. Como as respostas desses quatro modelos variam entre execuções, o Stuff foi executado três vezes com eles nas 25 perguntas de cada conjunto, para medir quanto da diferença entre métodos é só variação do próprio modelo. Os modelos da OpenAI foram acessados pela API da empresa, e os demais pelo OpenRouter \cite{openrouter2026}, com o provedor de cada chamada registrado; o Qwen rodou em precisão bf16 e o DeepSeek em fp8, a precisão em que é distribuído. No total, foram 12.150 execuções: 10.500 na camada principal (5 LLMs $\times$ 3 conjuntos $\times$ 7 condições $\times$ 100 perguntas), 1.050 na de generalização e 600 repetições do Stuff.

\subsection{Métricas}

\textbf{Métrica principal.} Cada conjunto tem uma métrica comparada diretamente com o gabarito, sem juiz automático:

\begin{itemize}
  \item \textbf{MedQA: acurácia.} A letra escolhida é extraída do início da resposta e comparada com a alternativa correta oficial.
  \item \textbf{FaQuAD: cobertura de tokens.} Fração dos tokens da resposta de referência que aparecem na resposta gerada, após a normalização do SQuAD (minúsculas, sem pontuação), vale a melhor entre as respostas aceitas. A cobertura corresponde à revocação das métricas oficiais do SQuAD \cite{rajpurkar2016squad}, que também foram calculadas: o \emph{exact match}, que vale 1 quando a resposta normalizada é idêntica a uma resposta aceita e 0 nos demais casos, e o F1, média harmônica entre a precisão (fração dos tokens da resposta gerada presentes na referência) e a revocação (fração dos tokens da referência presentes na resposta gerada). As duas, porém, penalizam respostas em frase completa diante de referências de poucas palavras: o \emph{exact match} do Stuff não passou de 0,12 em nenhum modelo, embora as respostas contivessem o trecho certo na maior parte dos casos. Os artigos não foram removidos na normalização, porque em português ``a'' e ``o'' também são preposição e pronome.
  \item \textbf{Docugami: valores do gabarito.} Fração dos valores numéricos distintos da resposta de referência --- montantes em dólares, percentuais e valores em milhões ou bilhões --- que aparecem na resposta gerada, aceitando arredondamento e troca de escala (US\$~7.442 milhões equivalem a US\$~7,44 bilhões), com tolerância de até 2\%. Datas, anos e numeração de itens não contam como valor.
\end{itemize}

A métrica do Docugami foi definida depois de se observar que a similaridade de cosseno dava notas altas a respostas genéricas, sem nenhum número: uma recusa parcial como ``\emph{I don't have verified figures}'', seguida de comentários gerais sobre a empresa, chegava a 0,86 de similaridade com uma referência cheia de valores. A \textbf{similaridade de cosseno}, métrica do estudo de referência, é o cosseno do ângulo entre o embedding da resposta gerada e o da resposta de referência: quanto mais próxima de 1, mais próximos estão os dois textos no espaço do embedding, o que reflete semelhança de tema e de vocabulário, e não a coincidência dos valores e fatos. Calculada com o BGE-M3, ela foi mantida em todas as execuções para comparação, e recalculada com o bge-small-en-v1.5, o embedding usado por aquele estudo, nos dois conjuntos em inglês.

\textbf{Recusas.} Uma resposta foi classificada como recusa quando o modelo declarou não ter a informação, total ou parcialmente (``\emph{I don't know}'', ``Não sei'', ``\emph{the context does not provide}\ldots'', ``\emph{I don't have verified figures}''), por expressões regulares tiradas das próprias respostas. A taxa de recusa é a fração das respostas classificadas como recusa.

\textbf{Custo.} Para cada resposta, foram registrados os tokens --- as unidades em que o tokenizador de cada modelo divide o texto, e pelas quais as APIs cobram --- de entrada (o prompt, com a instrução, a pergunta e os fragmentos) e de saída (o texto gerado), informados pela API, o custo em dólares, calculado com os preços de tabela da Tabela~\ref{tab:llms}, e o tempo total, da recuperação até o fim da geração. Nos métodos com várias chamadas, os tokens, o custo e o tempo de uma resposta somam todas as chamadas. O custo é nominal: não considera descontos de cache de prompt, que dependem do provedor e da ordem das chamadas. O uso de CPU e memória da máquina local não foi analisado, porque, com os LLMs acessados por API, essas medidas não refletem o custo de inferência dos métodos.

\textbf{Recuperação.} A qualidade da busca foi medida sem chamadas ao LLM, contra o trecho certo marcado no gabarito: o relatório citado na resposta de referência, no Docugami, e o parágrafo de origem da pergunta, no FaQuAD. Um fragmento recuperado conta como relevante quando pertence a esse relatório ou parágrafo. O Hit@4 é a fração das perguntas em que ao menos um fragmento relevante aparece entre os quatro primeiros recuperados, o mesmo $k$ usado pelos métodos. O Recall@4 é a fração dos documentos certos de cada pergunta que aparecem entre esses quatro fragmentos, em média sobre as perguntas; ele difere do Hit@4 no Docugami, em que 60 perguntas citam de dois a quatro relatórios, e coincide com ele no FaQuAD, em que cada pergunta tem um único parágrafo de origem. O MRR@10 (\emph{mean reciprocal rank}) é a média, sobre as perguntas, do inverso da posição do primeiro fragmento relevante entre os dez primeiros: vale 1 quando ele vem em primeiro lugar, $1/2$ quando vem em segundo e 0 quando não aparece, e indica em que posição o trecho certo costuma aparecer. As três métricas foram calculadas em todas as perguntas com o trecho marcado: 193 das 195 do Docugami, já que duas não citam o relatório de origem, e as 900 do FaQuAD. O MedQA não marca o trecho certo.

\textbf{RAGAS.} Foram calculadas quatro métricas do RAGAS \cite{es2024ragas} (biblioteca \texttt{ragas} 0.4.3), com o \texttt{gpt-4o-mini} como juiz e o \texttt{text-embedding-3-small} da OpenAI como embedding. A \emph{Faithfulness} mede se a resposta está sustentada pelos fragmentos: o juiz divide a resposta gerada em afirmações e verifica, para cada uma, se ela pode ser inferida dos fragmentos entregues ao LLM; a nota, de 0 a 1, é a fração sustentada, e valores baixos indicam conteúdo que não vem do contexto, como alucinações. Ela foi calculada nos quatro métodos que entregam ao LLM um conjunto único de fragmentos (Stuff, Refine, Map-Reduce e Map-Rerank); no Query Step-Down e no Reciprocal, cada pergunta alternativa recupera os seus próprios fragmentos, e não há um contexto único contra o qual verificar a resposta. A \emph{Answer Relevancy} mede se a resposta trata da pergunta feita: o juiz gera, a partir da resposta, três perguntas que ela responderia, e a nota é a média da similaridade de cosseno entre o embedding de cada uma e o da pergunta original; respostas que o juiz classifica como evasivas, como as recusas, recebem zero. Ela não verifica se a resposta está correta e foi calculada em todos os métodos e no controle sem RAG. \emph{Context Precision} e \emph{Context Recall} avaliam os fragmentos recuperados com a pergunta original, que são os mesmos nesses quatro métodos e em todos os LLMs; por isso foram calculadas uma única vez por pergunta, em vez de repetidas a cada linha, o que só mediria a variação do juiz. A \emph{Context Precision} mede se os fragmentos úteis vêm no topo da lista: o juiz decide, para cada um dos quatro fragmentos, se ele ajuda a chegar à resposta de referência, e a nota é a média da precisão nas posições dos fragmentos úteis, de modo que um fragmento útil em primeiro lugar pesa mais que um em quarto. A \emph{Context Recall} mede se os fragmentos bastam para responder: o juiz divide a resposta de referência em afirmações e verifica quais delas podem ser sustentadas pelos fragmentos recuperados, e a nota é a fração sustentada. As duas vão de 0 a 1 e dependem da resposta de referência, e não da resposta gerada. No MedQA, a referência dessas métricas é o texto da alternativa correta, e elas são a única avaliação da recuperação.

\textbf{Validação do juiz.} A \emph{Faithfulness} do juiz foi comparada com os rótulos de sustentação pelo contexto do RAGBench \cite{friel2024ragbench} em uma amostra estratificada de 206 respostas de sete subconjuntos, com cerca de metade de respostas sustentadas e metade com alucinação em cada um. A concordância foi medida pela área sob a curva ROC (AUROC) \cite{hanley1982roc}, a probabilidade de que, sorteadas uma resposta sustentada e uma com alucinação, a sustentada receba a maior nota: 0,5 equivale ao acaso, e 1, à separação perfeita. A AUROC foi de 0,701, contra 0,567 da \emph{Faithfulness} calculada pelos próprios autores do RAGBench sobre as mesmas respostas, e variou de 0,54 a 0,90 entre os subconjuntos. O juiz é, portanto, moderadamente confiável, e as métricas RAGAS entram como complementares: as conclusões principais vêm das métricas comparadas com o gabarito. Como o \texttt{gpt-4o-mini} também é um dos modelos avaliados, as suas notas RAGAS não foram usadas para comparar LLMs, por risco de autopreferência \cite{panickssery2024selfpref}.

\subsection{Bancos vetoriais}

Como o índice HNSW do ChromaDB é aproximado, a comparação com o FAISS (\texttt{IndexFlatL2}, busca exata) mede tanto a infraestrutura quanto o efeito da aproximação sobre os resultados. Os dois bancos indexaram os mesmos vetores BGE-M3 dos três conjuntos, que formam três escalas de base (318, 4.463 e 121.553 vetores), em processos separados e sem nada mais em execução, e responderam às consultas das perguntas da amostra. Foram medidos o tempo de indexação, o tempo médio por consulta, a memória alocada na indexação (o aumento da memória residente do processo) e a concordância do top-4 entre os bancos, isto é, a fração das consultas em que os dois devolvem os mesmos quatro fragmentos; no Docugami e no FaQuAD, também o Hit@4 de cada banco. Por fim, nas perguntas em que o top-4 do FAISS diferiu do recebido nos experimentos, o Stuff do \texttt{gpt-4o-mini} foi executado de novo com os fragmentos do FAISS, para medir o efeito da troca de banco na métrica principal.

\subsection{Análise estatística}

Em cada combinação de conjunto e LLM, as perguntas foram tratadas como blocos pareados, já que todos os métodos respondem às mesmas perguntas. A diferença entre os seis métodos foi testada com o teste de Friedman \cite{friedman1937ranks}, que ordena os métodos dentro de cada pergunta e verifica se os postos médios diferem mais do que o acaso explicaria, e, em seguida, com testes de Wilcoxon pareados \cite{wilcoxon1945individual} par a par, que comparam duas condições pelos postos das diferenças entre elas em cada pergunta, sem supor distribuição normal. Os valores-$p$ foram corrigidos pelo método de Holm \cite{holm1979simple}, que os ajusta em sequência para manter em no máximo 5\% a probabilidade de ao menos um falso positivo entre as comparações de uma mesma família ($p_{\text{Holm}}$ nas tabelas), conforme a recomendação de Demšar~\cite{demsar2006statistical}. O efeito do RAG foi testado da mesma forma, comparando cada método com o controle sem RAG, com a correção de Holm sobre as seis comparações. O posto médio de um método é a média da sua posição entre os seis (1 = melhor) nos LLMs, e a concordância dessa ordem entre os LLMs foi medida pelo coeficiente W de Kendall \cite{kendall1939mrankings}, de 0 (nenhuma concordância) a 1 (mesma ordem em todos), com o teste pela aproximação qui-quadrado. Comparações pareadas de acurácia entre duas condições usaram o teste exato de McNemar \cite{mcnemar1947note}, que considera só as perguntas em que uma condição acerta e a outra erra. A associação entre duas métricas foi medida pela correlação de Spearman \cite{spearman1904proof} ($\rho$), a correlação entre os postos das duas, de $-1$ a 1, em que 0 indica ausência de relação monotônica. Intervalos de confiança de 95\% foram estimados por \emph{bootstrap} percentil \cite{efron1993bootstrap} com 10.000 reamostragens. O valor-$p$ é a probabilidade de observar uma diferença ao menos tão grande quanto a medida se as condições fossem equivalentes, e o nível de significância adotado foi 0,05. Na camada de generalização, com 25 perguntas por conjunto (14 no Docugami), só diferenças grandes podem ser detectadas.

\subsection{Ambiente e reprodutibilidade}

Os experimentos foram executados entre outubro e novembro de 2026, em Python 3.12, em um notebook com processador Intel Core i5-12450H (12 \emph{threads}), cerca de 20 GB de RAM e sem GPU dedicada. Os embeddings dos dois corpora maiores foram calculados em GPU no Google Colab e conferidos contra os fragmentos locais por um \emph{hash} dos textos. As principais bibliotecas foram \texttt{chromadb} 1.5.9, \texttt{faiss} 1.14.3, \texttt{FlagEmbedding} 1.4.0, \texttt{openai} 2.41.1, \texttt{ragas} 0.4.3 e \texttt{scipy} 1.17.1. Cada execução foi registrada em um arquivo JSONL, com a pergunta, os fragmentos recuperados, a resposta, os tokens, os tempos, a versão do modelo informada pela API e o provedor. O custo nominal da geração foi de US\$~20,09. O código, os registros e os scripts que geram todas as tabelas e figuras deste artigo estão no repositório do trabalho.
